% kernel/engineering_knowledge.tex

\chapter{Engineering Knowledge}
\label{chap:engineering-knowledge}

Engineering activity requires information that may be relied upon for
interpretation, evaluation, design, and decision-making. The Knowledge
Kernel distinguishes such accepted information from observations,
representations, and proposals.

\section{Engineering Knowledge}

\begin{definition}[Engineering knowledge]
\label{def:engineering-knowledge}
Engineering knowledge is accepted engineering information that may be
used to interpret, evaluate, constrain, or decide engineering
questions within its applicable context.
\end{definition}

Engineering knowledge may originate from validated observations,
approved standards, established theory, accepted derivations,
regulatory rules, operational experience, or prior engineering
decisions.

\section{Knowledge and Evidence}

Evidence supports knowledge but is not identical to it. Observations,
calculations, and representations may contribute to an engineering
conclusion while retaining their individual provenance and uncertainty.

Acceptance establishes that information may be used with a specified
engineering authority and scope.

\section{Knowledge Ownership}

\begin{definition}[Knowledge ownership]
\label{def:knowledge-ownership}
Knowledge ownership is the responsibility assigning engineering
knowledge to the concept or authority that is entitled to establish,
maintain, and revise it.
\end{definition}

Ownership prevents engineering rules, measured reality, operational
assumptions, and solver-specific artefacts from being merged into an
undifferentiated information store.

\section{Truth Mode}

Engineering information may be observed, derived, assumed, proposed,
accepted, or decided. Its truth mode must remain explicit because these
statuses carry different authority.

A proposal is not yet a decision. A derived quantity is not an
observation. A representation is not knowledge merely because it
contains values.

\section{Canonical Interpretation}

\begin{proposition}
Engineering knowledge is distinguished by accepted authority and
applicability, not by storage format or computational origin.
\end{proposition}

\begin{proposition}
Engineering observations and representations may support knowledge but
do not replace the process by which knowledge is accepted.
\end{proposition}

\begin{proposition}
Knowledge ownership determines responsibility for the establishment and
revision of engineering knowledge.
\end{proposition}

Engineering knowledge provides the authoritative basis on which
constraints are evaluated, proposals are assessed, and engineering
decisions are made.
